\origpage{}{}
\begin{quote}
{\centering\textbf{Abstract}\par}

A linkage $L$ in a graph $G$ is a subgraph each component of which is a path, and it is \textit{vital} if $V(L) = V(G)$ and there is no other linkage in $G$ joining the same pairs of vertices. We show that, if $G$ has a vital linkage with $p$ components, then $G$ has tree-width bounded above by a function of $p$. This is the major step in the proof of the unproved lemma from Graph Minors XIII, and it has a number of other applications, including a constructive proof of the intertwining conjecture.
\end{quote}
\clearpage
